\documentclass[journal,10pt]{IEEEtran}
\setlength{\textwidth}{182mm}
\setlength{\columnsep}{4.2mm}
\usepackage{amsmath,amssymb,newtxmath}
\usepackage[no-math]{fontspec}
\setmainfont{Times New Roman}
\usepackage{xeCJK}
\setCJKmainfont[AutoFakeBold=2]{SimSun}
\usepackage{graphicx,booktabs,array,hyperref,placeins}
\hypersetup{hidelinks,unicode=true}
\raggedbottom
\setlength{\emergencystretch}{2em}
\renewcommand{\arraystretch}{1.15}
\usepackage{seqsplit}
\begin{document}
\title{Hierarchical Verification, Common-Terminal Current Sharing, and Finite-Length Longitudinal PEEC for Litz-Wire Modeling · S1 Theory and Discretization}
\author{Miao Yuxuan\thanks{Huainan Normal University}}

\maketitle
\noindent 合并稿对应章节 / Chapter in merged manuscript: \href{paper_en_ieee.pdf#chapter-S1}{S1}\par
\FloatBarrier
\section*{S1 Theory and Discretization}
\hypertarget{chapter-S1}{}
\pdfbookmark[0]{S1 Theory and Discretization}{bookmark-chapter-S1}
\subsection*{Round Derivation}
\hypertarget{S1.round_derivation.heading}{}
\pdfbookmark[1]{Round Derivation}{bookmark-S1.round_derivation.heading}
The axial current in a round wire satisfies (1/r)d(r dJz/dr)/dr=κ²Jz. Regularity at the axis selects I\(_0\)(κr). The identity \(\int\)\(_0\)\(^a\) rI\(_0\)(κr)dr=aI\(_1\)(κa)/κ together with the total-current constraint determines the normalization. The surface axial electric field divided by total current gives the internal impedance; its real part is checked by an independent Joule integral. Thus the derivation connects analytical shape, normalization, and power.


\subsection*{Sharing Derivation}
\hypertarget{S1.sharing_derivation.heading}{}
\pdfbookmark[1]{Sharing Derivation}{bookmark-S1.sharing_derivation.heading}
Sufficiency of the sharing criterion follows from uniqueness of the port solution. If the projection condition holds, every component of (It/N)Z1+e is equal, so a common voltage exists for which the equal-current vector satisfies both the port equations and the total-current constraint. Necessity follows by substitution of an existing equal-current solution. Numerical sharing is assessed using normalized complex-current deviation, with magnitude and phase reported to retain the information in phasor addition.


\subsection*{Dtm Derivation}
\hypertarget{S1.dtm_derivation.heading}{}
\pdfbookmark[1]{Dtm Derivation}{bookmark-S1.dtm_derivation.heading}
Let Ps map fixed physical strand identities to spatial positions in section s. The local current is Ps i and the voltage in identity coordinates is Ps\(^{\mathrm T}\)ZsPs i. Summation over series segments Δzs yields the total impedance. Position permutation and current identity therefore have separate records; K, 2K, and 4K segmentation checks concern the same paths and physical strands. Period, gap, and arc-length checks constrain geometric feasibility and copper usage separately.


\subsection*{Peec Derivation}
\hypertarget{S1.peec_derivation.heading}{}
\pdfbookmark[1]{Peec Derivation}{bookmark-S1.peec_derivation.heading}
Within each volume filament, current is represented by the constant axial density iα/Aα. Projection of the volume-integral vector potential onto the testing filament and integration along length produces the symmetric partial-inductance matrix. The self kernel is integrable but singular at coincidence and uses a dedicated accuracy setting; mutual terms use an independent order. Combining finite-segment ohmic resistance with jωL forms Z, after which the common-terminal augmented system determines filament currents and terminal voltage.


\subsection*{Triangle Derivation}
\hypertarget{S1.triangle_derivation.heading}{}
\pdfbookmark[1]{Triangle Derivation}{bookmark-S1.triangle_derivation.heading}
On a linear triangle, J=\(\sum\)\(_a\)NaJa. Current integration uses A/3 for each shape function. Joule integration uses the mass matrix and Joule expression in \href{paper_en_ieee.pdf#E09}{(8)}. Material markers assign each cell to a strand, the return, or air. Coordinates are converted to metres before area evaluation, and the integrated power is independently compared with the archived scalar.


\hypertarget{E06}{}\begin{equation}Z_{\mathrm{tot}}=\sum_{s=1}^{K}\Delta z_s P_s^T Z_s P_s\tag{7}\label{eq:E06}\end{equation}
\hypertarget{E09}{}
\begin{equation}
\begin{gathered}
M_T=\frac{A_T}{12}\begin{bmatrix}2&1&1\\1&2&1\\1&1&2\end{bmatrix} \\
P_T=\mathbf J_T^H M_T\mathbf J_T/\sigma
\end{gathered}
\tag{8}\label{eq:E09}
\end{equation}
\end{document}
